\documentclass[10pt,a4paper]{amsart}
\usepackage[margin=19mm]{geometry}
\usepackage{amsmath,amssymb,mathtools}
\usepackage[scr=rsfs,cal=euler]{mathalfa}
\usepackage{xfrac}
\usepackage[unicode,hidelinks,bookmarks=false]{hyperref}
\newcommand{\ie}{i.e.\ }
\newcommand{\cf}{cf.\ }
\newcommand{\resp}{respectively\ }
\newcommand{\Subsection}{\subsection}
\providecommand{\nobreakdash}{\nobreak\hbox{-}}
\setlength{\emergencystretch}{2em}
\multlinegap0pt
\usepackage{eufrak}
\usepackage{amscd}
\usepackage{amstext}
\usepackage{xy}
 \newcommand{\Q}{{\mathcal Q}}
 \newcommand{\K}{{\mathbb K}}
 \renewcommand{\L}{{\mathbb L}}
   \newcommand{\Sp}{\operatorname{Sp}}
   \newcommand{\Spec}{\operatorname{Spec}}
   \newcommand{\gr}{\operatorname{gr}}
   \newcommand{\Hom}{\operatorname{Hom}}
   \newcommand{\CH}{\operatorname{CH}}
 \newcommand{\Ext}{\operatorname{Ext}}
\newcommand{\Ad}{\operatorname{Ad}}
\newcommand{\eq}[1]{\begin{equation}#1\end{equation}}
\newcommand{\eQ}[1]{\begin{equation*}#1\end{equation*}}
\newcommand{\spl}[1]{\begin{split}#1\end{split}}
 \newcommand{\Sup}{\operatorname{Sup}}
 \newcommand{\Max}{\operatorname{Max}}
\newcommand{\Ker}{\operatorname{Ker}}
\newcommand{\id}{\operatorname{id}}
\newcommand{\cone}{\operatorname{cone}}
\newcommand{\im}{\operatorname{im}}
\newcommand{\ev}{\operatorname{ev}}
\newcommand{\HS}{\operatorname{HS}}
\newcommand{\Tr}{\operatorname{Tr}}
\newcommand{\var}{\operatorname{var}}
\newcommand{\supp}{\operatorname{supp}}
\newcommand{\diam}{\operatorname{diam}}
\newcommand{\nnn}[1]{|\! |\! | #1 |\! |\! |}
   \newtheorem{thm}{Theorem}[section]
   \newtheorem{prop}[thm]{Proposition}
   \newtheorem{lem}[thm]{Lemma}
   \newtheorem{cor}[thm]{Corollary}
   \newtheorem{exerc}[thm]{Exercise}
   \newtheorem{defn}[thm]{Definition}
   \newtheorem{example}[thm]{Example}
   \newtheorem{obs}[thm]{Observation}
   \newtheorem{remark}[thm]{Remark}
	\newtheorem{problem}[thm]{Problem}
 \numberwithin{equation}{section}
\begin{document}
\title[On ideals in the semilattice of coarse equivalence classes o]{On ideals in the semilattice of coarse equivalence classes of metrics}
\author{Vladimir Manuilov}
\date{}
\maketitle
\noindent\emph{Source and licence.} On ideals in the semilattice of coarse equivalence classes of metrics. Version arXiv:2606.27816v1 (CC BY 4.0). This material is distributed under the Creative Commons Attribution 4.0 International licence, http://creativecommons.org/licenses/by/4.0/, as recorded in the arXiv OAI licence metadata for this exact version and shown on the abstract page https://arxiv.org/abs/2606.27816v1. Full rights notice: licenses/arxiv-2606.27816-CC-BY-4.0.txt.\par\smallskip
\noindent\textit{Editorial scope.} Counted unit: the original \texttt{thm} environment \texttt{-} at source lines 525--529 together with its complete proof at lines 530--550; the delivered document keeps source lines 519--551 so that every referenced label and every citation resolves inside it. Native editable source is the paper's own concatenated TeX; the AMS wrapper is editorial formatting only. No publisher class, upstream script or unrelated artwork is redistributed.
\section{Bounded geometry ideals} 

Recall that a metric $a$ on $X$ is of \emph{bounded geometry} if for any $R>0$ there exists $C>0$ such that every ball of radius R contains no more than $C$ points.

Let $F^{BG}\subset M(X)$ denote the subset of coarse equivalence classes of metrics of bounded geometry, and let $F^{BG(\alpha)}=F^{BG}\cap F_\alpha$, $\alpha\in M(X)$.


\begin{thm}
$F^{BG(\alpha)}$ is an ideal in $M(X)$ for any $\alpha\in M(X)$. Moreover, $F^{BG(\alpha)}$ lies in the closure of principal ideals, i.e. $F^{BG(\alpha)}\in\overline{M(X)}\subset I(M(X))$.
The ideal $F^{BG(\alpha)}$ is principal iff $\alpha$ is of bounded geometry.

\end{thm}

\begin{proof}
Bounded geometry is preserved under coarse equivalence. If $b\preceq c$ in $\mathcal M(X)$ and if $c$ has bounded geometry then $b$ has bounded geometry too. It was proved in \cite{Man-Zhu}, Prop. ??? that if $b$ and $c$ have bounded geometry then $b\lor c$ has bounded geometry too. Therefore, $F^{BG}$ is an ideal. Then $F^{BG(\alpha)}$ is an ideal as well. 


Let us show that $F^{BG}$ lies in the closure of principal ideals. As a base of neighborhoods of $F\in I(M(X))$ one can take the sets 
$$
U_{\alpha_1,\ldots,\alpha_n;\beta_1,\ldots,\beta_m}(F)=\{F'\in I(M(X)): \alpha_i\in F', \beta_j\notin F'\}, 
$$  
where $\alpha_i,\beta_j\in M(X)$, $\alpha_i\in F$, $\beta_j\notin F$, $i=1,\ldots,n$, $j=1,\ldots,m$. 
Let $\alpha_i$, $i=1,\ldots,n$, have bounded geometry, and $\beta_j$, $j=1,\ldots,m$, do not have bounded geometry. Set $\gamma=\alpha_1\lor\cdots\lor \alpha_n$. Then $\gamma$ has bounded geometry, and $\alpha_i\in F_\gamma$ for any $i=1,\ldots,n$. If $\beta\in F_\gamma$ then $\beta\preceq \gamma$, hence $\beta$ is of bounded geometry. Therefore, $\beta_j\notin F_\gamma$, $j=1,\ldots,m$. Thus, $F_\gamma\in U_{\alpha_1,\ldots,\alpha_n;\beta_1,\ldots,\beta_m}(F^{BG})$. Hence there exists a net $\{F_{\gamma_\lambda}\}$ that converges to $F^{BG}$. As the map $F\to F\cap G$ is continuous on $I(M(X))$ for any $G\in I(M(X))$, the net $\{F_{\gamma_\lambda}\cap F_\alpha\}$ converges to $F^{BG(\alpha)}$.

If $\alpha$ is of bounded geometry then any (equivalence class of) metric $b$ with $\beta=[b]\preceq \alpha$ is of bounded geometry, hence $F_\alpha$ consists of metrics of bounded geometry, i.e. coincides with $F^{BG(\alpha)}$. 

Now let $\alpha$ be not of bounded geometry, and let $a\in\mathcal M(X)$ be a representative for $\alpha\in M(X)$. We have to show that $F^{BG(\alpha)}$ is not principal. Suppose the contrary: $F^{BG(\alpha)}=F_\beta$ for some $\beta\in M(X)$. As $\beta\in F_b$, $\beta$ has bounded geometry, and $\beta\preceq \alpha$. Let $b\in\mathcal M(X)$ be a representative for $\beta$. To get a contradiction, we want to find a bounded geometry metric $c\in\mathcal M(X)$ such that $\gamma=[c]$ satisfies $\gamma\preceq \alpha$ (then $\gamma\in F^{BG(\alpha)}$) and $\gamma\npreceq \beta$ (then $\gamma\notin F_\beta$). 

As $\beta\preceq \alpha$ and $\beta\neq \alpha$ (otherwise they both would be of unbounded geometry), there exists $R>0$ and two sequences $\{x_n\}$ and $\{y_n\}$, $n\in\mathbb N$, of points in $X$ such that $a(x_n,y_n)\leq R$ and $b(x_n,y_n)>n$ for any $n\in\mathbb N$. Set 
$$
b'(x,y)=\left\lbrace\begin{array}{cl}0,&\mbox{if\ }y=x;\\R,&\mbox{if\ }\{x,y\}=\{x_n,y_n\}\mbox{\ for\ some\ }n\in\mathbb N;\\\infty,&\mbox{otherwise}.\end{array}\right.
$$   
This is only a pseudometric, not a metric, but $b\lor b'$ is a well defined metric. Set $c=b\lor b'$. As $b'(x,y)\geq a(x,y)$ for any $x,y\in X$, $c\preceq a$. As $c(x_n,y_n)\leq R$, while $b(x_n,y_n)>n$, $n\in\mathbb N$, we have $c\npreceq b$.    
\end{proof}


\begin{thebibliography}{99}
\bibitem[Man-Zhu]{Man-Zhu} V. Manuilov, J. Zhu. \textit{Two versions of the uniform Roe algebras for spaces without bounded geometry.} Math. Nachr. 294 (2021), 1140--1147.
\end{thebibliography}
\end{document}
